% GENERATED FILE. Edit canonical lessons, then run npm run generate:books.
% canonical-source-hash: fnv1a64:9b65e3fd03e1c9c1
% canonical-lessons: LA-C354-assus, LA-C354-lucanica, LA-C354-craticula, LA-C354-frigo, LA-C354-torreo, LA-R354-review-347-354

\chapter{Assum et frīctum}
\label{ch:la-assum-et-frictum}
\hlchaptermodality{354}

\begin{chapteropening}
\textbf{By the end of this chapter:} \emph{I can describe roasting and frying: roasted, a sausage, the gridiron, to fry and to toast.}

Complete the last lesson of chapter 354: I can describe roasting and frying: roasted, a sausage, the gridiron, to fry and to toast.
\end{chapteropening}

\section[assus, assa, assum]{assus, assa, assum — roasted}

\label{lesson:LA-C354-assus}



\begin{quote}
Before the new one: say the Latin for foam, then the Latin for to skim off.
\end{quote}

\subsection*{You'll want to know: assus, assa, assum}
\textbf{assus, assa, assum} — \textquotedblleft{}roasted\textquotedblright{}.

\begin{quote}
Carō assa iūcunda erat.
\end{quote}

\begin{quote}
\emph{The roast meat was delicious.}
\end{quote}

\emph{Assus} is roasted, \emph{ēlixus} boiled: Horace and Celsus set the two side by side.

\begin{cousinweb}
\textbf{assus} is inherited (EtymDB); Lewis \& Short connect it with \emph{āreō}, \textquotedblleft{}to be dry\textquotedblright{}, and \emph{ārdeō}, \textquotedblleft{}to burn\textquotedblright{}.
\end{cousinweb}

\begin{grammarlens}[title={What you've built}]
One more word for telling, explaining and describing at B1.
\end{grammarlens}

\subsection*{Your turn}
\begin{itemize}
\raggedright
  \item \emph{Say it:} \emph{assus, assa, assum}
  \item \emph{Say it:} \emph{assus, assa, assum}, once more
  \item \emph{Say it:} a sentence of your own with \emph{assus, assa, assum}
  \item \emph{Recall:} say the Latin for foam, then the Latin for to skim off, then say \emph{assus, assa, assum} again
\end{itemize}

\begin{tcolorbox}[breakable,colback=teal!4,colframe=teal!35!black,title={Before you move on}]
What does assus, assa, assum mean? (\textquotedblleft{}Roasted\textquotedblright{}.) Say it once more.
\end{tcolorbox}
\section[lūcānica, lūcānicae]{lūcānica, lūcānicae — a smoked sausage}

\label{lesson:LA-C354-lucanica}



\begin{quote}
Before the new one: say the Latin for to skim off, then the Latin for roasted.
\end{quote}

\subsection*{You'll want to know: lūcānica, lūcānicae}
\textbf{lūcānica, lūcānicae} — \textquotedblleft{}a smoked sausage\textquotedblright{}.

\begin{quote}
Senex lūcānicam et pānem edēbat.
\end{quote}

\begin{quote}
\emph{The old man was eating a sausage and bread.}
\end{quote}

Feminine. A smoked pork sausage from southern Italy; Martial sends one as a little present.

\begin{cousinweb}
\textbf{lūcānica} is the sausage of Lucania, a region of southern Italy, named after the people who made it (EtymDB).
\end{cousinweb}

\begin{grammarlens}[title={What you've built}]
One more word for telling, explaining and describing at B1.
\end{grammarlens}

\subsection*{Your turn}
\begin{itemize}
\raggedright
  \item \emph{Say it:} \emph{lūcānica, lūcānicae}
  \item \emph{Say it:} \emph{lūcānica, lūcānicae}, once more
  \item \emph{Say it:} a sentence of your own with \emph{lūcānica, lūcānicae}
  \item \emph{Recall:} say the Latin for to skim off, then the Latin for roasted, then say \emph{lūcānica, lūcānicae} again
\end{itemize}

\begin{tcolorbox}[breakable,colback=teal!4,colframe=teal!35!black,title={Before you move on}]
What does lūcānica, lūcānicae mean? (\textquotedblleft{}A smoked sausage\textquotedblright{}.) Say it once more.
\end{tcolorbox}
\section[crāticula, crāticulae]{crāticula, crāticulae — a gridiron, a grill}

\label{lesson:LA-C354-craticula}



\begin{quote}
Before the new one: say the Latin for roasted, then the Latin for a smoked sausage.
\end{quote}

\subsection*{You'll want to know: crāticula, crāticulae}
\textbf{crāticula, crāticulae} — \textquotedblleft{}a gridiron, a grill\textquotedblright{}.

\begin{quote}
Piscēs in crāticulā super carbōnēs coquēbantur.
\end{quote}

\begin{quote}
\emph{The fish were cooking on a grill over the coals.}
\end{quote}

Feminine, a small grid of iron. Martial makes a small present of one.

\begin{cousinweb}
\textbf{crāticula} is a diminutive of \emph{crātis}, \textquotedblleft{}wickerwork, a hurdle\textquotedblright{} (Lewis \& Short): a little grid.
\end{cousinweb}

\begin{grammarlens}[title={What you've built}]
One more word for telling, explaining and describing at B1.
\end{grammarlens}

\subsection*{Your turn}
\begin{itemize}
\raggedright
  \item \emph{Say it:} \emph{crāticula, crāticulae}
  \item \emph{Say it:} \emph{crāticula, crāticulae}, once more
  \item \emph{Say it:} a sentence of your own with \emph{crāticula, crāticulae}
  \item \emph{Recall:} say the Latin for roasted, then the Latin for a smoked sausage, then say \emph{crāticula, crāticulae} again
\end{itemize}

\begin{tcolorbox}[breakable,colback=teal!4,colframe=teal!35!black,title={Before you move on}]
What does crāticula, crāticulae mean? (\textquotedblleft{}A gridiron, a grill\textquotedblright{}.) Say it once more.
\end{tcolorbox}
\section[frīgō, frīgere]{frīgō, frīgere — to fry, to parch}

\label{lesson:LA-C354-frigo}



\begin{quote}
Before the new one: say the Latin for a smoked sausage, then the Latin for a gridiron.
\end{quote}

\subsection*{You'll want to know: frīgō, frīgere}
\textbf{frīgō, frīgere} — \textquotedblleft{}to fry, to parch\textquotedblright{}.

\begin{quote}
Coquus piscēs in oleō frīxit.
\end{quote}

\begin{quote}
\emph{The cook fried the fish in oil.}
\end{quote}

\emph{frīgō, frīgere, frīxī, frīctum}, long \emph{ī}, in a \emph{sartāgō} (taught). English \emph{fry} comes from it through French.

\begin{cousinweb}
\textbf{frīgō} is compared by Lewis \& Short with Greek \emph{phrygō}, \textquotedblleft{}to roast, to parch\textquotedblright{}. English \emph{fry} comes from it through French.
\end{cousinweb}

\begin{grammarlens}[title={What you've built}]
One more word for telling, explaining and describing at B1.
\end{grammarlens}

\subsection*{Your turn}
\begin{itemize}
\raggedright
  \item \emph{Say it:} \emph{frīgō, frīgere}
  \item \emph{Say it:} \emph{frīgō, frīgere}, once more
  \item \emph{Say it:} a sentence of your own with \emph{frīgō, frīgere}
  \item \emph{Recall:} say the Latin for a smoked sausage, then the Latin for a gridiron, then say \emph{frīgō, frīgere} again
\end{itemize}

\begin{tcolorbox}[breakable,colback=teal!4,colframe=teal!35!black,title={Before you move on}]
What does frīgō, frīgere mean? (\textquotedblleft{}To fry, to parch\textquotedblright{}.) Say it once more.
\end{tcolorbox}
\section[torreō, torrēre]{torreō, torrēre — to roast, to toast, to parch}

\label{lesson:LA-C354-torreo}



\begin{quote}
Before the new one: say the Latin for a gridiron, then the Latin for to fry.
\end{quote}

\subsection*{You'll want to know: torreō, torrēre}
\textbf{torreō, torrēre} — \textquotedblleft{}to roast, to toast, to parch\textquotedblright{}.

\begin{quote}
Ad focum pānem torrēbat.
\end{quote}

\begin{quote}
\emph{He was toasting bread by the hearth.}
\end{quote}

\emph{torreō, torrēre, torruī, tostum}. English \emph{toast} comes from \emph{tostum}, and \emph{torrid} from the verb.

\begin{cousinweb}
\textbf{torreō} is inherited from Proto-Italic \emph{*torzeō}, \textquotedblleft{}to scorch\textquotedblright{} (EtymDB), a cousin of English \emph{thirst}. English \emph{toast} and \emph{torrid} come from it.
\end{cousinweb}

\begin{grammarlens}[title={What you've built}]
One more word for telling, explaining and describing at B1.
\end{grammarlens}

\subsection*{Your turn}
\begin{itemize}
\raggedright
  \item \emph{Say it:} \emph{torreō, torrēre}
  \item \emph{Say it:} \emph{torreō, torrēre}, once more
  \item \emph{Say it:} a sentence of your own with \emph{torreō, torrēre}
  \item \emph{Recall:} say the Latin for a gridiron, then the Latin for to fry, then say \emph{torreō, torrēre} again
\end{itemize}

\begin{tcolorbox}[breakable,colback=teal!4,colframe=teal!35!black,title={Before you move on}]
What does torreō, torrēre mean? (\textquotedblleft{}To roast, to toast, to parch\textquotedblright{}.) Say it once more.
\end{tcolorbox}
\section[(dialogue)]{Words from chapters 347–354 — a review}

\label{lesson:LA-R354-review-347-354}



\begin{quote}
Say the Latin for a monster and a giant.
\end{quote}

\subsection*{Your turn}
Say these aloud:

\begin{itemize}
\raggedright
  \item a monster, a giant, a hero, a nymph, a fleece
  \item an oracle, a seer, a portent, an omen, to prophesy
  \item to chat, to say yes, wonderful, unbelievable, an invention
  \item to buy food, food to go with bread, to wash off, to wipe, a small knife
  \item to chop up, to cut, a mortar, a pestle, to grind
  \item a millstone, a sieve, to knead, leaven, a baking pot
  \item to boil, a cooking pot, boiled, foam, to skim off
  \item roasted, a smoked sausage, a gridiron, to fry, to roast
\end{itemize}

\begin{tcolorbox}[breakable,colback=teal!4,colframe=teal!35!black,title={Before you move on}]
A monster? (\textbf{mōnstrum, mōnstrī}.) A giant? (\textbf{gigās, gigantis}.) A hero? (\textbf{hērōs, hērōis}.) A nymph? (\textbf{nympha, nymphae}.) A fleece? (\textbf{vellus, velleris}.) An oracle? (\textbf{ōrāculum, ōrāculī}.) A seer? (\textbf{vātēs, vātis}.) A portent? (\textbf{prōdigium, prōdigiī}.) An omen? (\textbf{ōmen, ōminis}.) To prophesy? (\textbf{vāticinor, vāticinārī}.) To chat? (\textbf{fābulor, fābulārī}.) To say yes? (\textbf{āiō}.) Wonderful? (\textbf{mīrābilis, mīrābile}.) Unbelievable? (\textbf{incrēdibilis, incrēdibile}.) An invention? (\textbf{commentum, commentī}.) To buy food? (\textbf{obsōnō, obsōnāre}.) Food to go with bread? (\textbf{obsōnium, obsōniī}.) To wash off? (\textbf{abluō, abluere}.) To wipe? (\textbf{tergeō, tergēre}.) A small knife? (\textbf{cultellus, cultellī}.) To chop up? (\textbf{concīdō, concīdere}.) To cut? (\textbf{secō, secāre}.) A mortar? (\textbf{mortārium, mortāriī}.) A pestle? (\textbf{pistillum, pistillī}.) To grind? (\textbf{terō, terere}.) A millstone? (\textbf{mola, molae}.) A sieve? (\textbf{crībrum, crībrī}.) To knead? (\textbf{subigō, subigere}.) Leaven? (\textbf{fermentum, fermentī}.) A baking pot? (\textbf{clībanus, clībanī}.) To boil? (\textbf{ferveō, fervēre}.) A cooking pot? (\textbf{caccabus, caccabī}.) Boiled? (\textbf{ēlixus, ēlixa, ēlixum}.) Foam? (\textbf{spūma, spūmae}.) To skim off? (\textbf{dēspūmō, dēspūmāre}.) Roasted? (\textbf{assus, assa, assum}.) A smoked sausage? (\textbf{lūcānica, lūcānicae}.) A gridiron? (\textbf{crāticula, crāticulae}.) To fry? (\textbf{frīgō, frīgere}.) To roast? (\textbf{torreō, torrēre}.)
\end{tcolorbox}
